\chapter{Conclusions}
\label{cha:conclusions}

Current cluster architectures are mainly concerned with maximizing service performance
and achieving near-instantaneous responsiveness while ignoring the energy impact
and the various resources involved. As a result, it is becoming more and more necessary---almost
a requirement---to reconsider and transform existing designs to create architectures
that are more resource-aware and sustainable. These latter requirements must be
considered just as vital as performance and responsiveness. \\ %
The context and fundamental principles on which the new and more sustainable
cluster architectures should either completely or partially rely on have been
initially explained in this document. Moreover, due to the outsourcing
phenomenon that shifts local and physical architectures towards remote, intangible,
and virtual architectures, users tend to forget the numerous implications of this
decision that not only can have privacy and/or security concerns but also, and
most importantly in this case, completely removes the direct control and the resource
awareness of the overall system. \\ %
Following the introduction, a possible theoretical design that is more
sustainable has been put forward, prioritizing the reduction of resource waste and
energy consumption while maintaining a certain degree of adaptability and configurability.
Every component that constitutes the architecture has been thoroughly explained,
highlighting its main purpose and different capabilities. A cluster is not a
cluster if the different components are not connected. In the design, there are
three major network groups; the External Network may even be removed if external
resources and connections are not required. The reCluster project, which is a practical
and sustainable cluster implementation of the theorized architecture, has also
been presented. All of its hardware components have been recycled following decommissioning.
\\ %
The implementation thoroughly explains how the reCluster was put into practice while
adhering to both the fundamental requirements and the theoretical architecture behind
the networks and components. It begins with the foundations, which are composed of
GNU/Linux, a collection of programs, and the Init System, all of which are
packaged together as a custom distribution into a portable and simple-to-install
ISO image file. Then, each theorized component is mapped to a real-world
application, and some of its most important features used in reCluster are explained.
The way the Server and Cluster Autoscaler components are implemented is
different because the first one was entirely built from scratch, whereas the second
one used an existing core that had been modified to be compatible with the cluster
version. The Server component represents the low-level knowledge of the cluster
and is deeply described, including the database structure, GraphQL queries and
mutations API, how the upscaling and downscaling procedures are handled, and how
to exploit Kubernetes monitoring without the need for a custom implementation. The
Cluster Autoscaler, on the other hand, represents the high-level knowledge of the
cluster, and it is used in conjunction with the Server component to perform the
upscaling or downscaling operations. The Cluster Autoscaler only scales the number
of nodes and, because it is not the only method of autoscaling in a cluster, the
other two approaches are also explained: Vertical Pod Autoscaler, which automatically
adjusts the resources of a Pod, and Horizontal Pod Autoscaler, which automatically
adjusts the number of Pods. The combination of the Horizontal Pod Autoscaler and
the Cluster Autoscaler is critical for improving total automation and minimizing
the need for human intervention, as well as lowering overall power consumption while
maintaining scalability. \\ %
Lastly, because manually installing the cluster and providing all of the
necessary information is tedious, the installation script that automatically performs
the installation operation while following a specific set of configuration parameters
has been discussed and proposed. Moreover, it is demonstrated how the node's benchmarks
and power consumption data are measured. Finally, the most essential elements of
the configuration and deployment files of the various components and/or technologies
used are illustrated and discussed, so that the employed features of each can be
recognized. \\ %
With all of the prior information and technology, building and deploying a more sustainable
and resource-aware cluster is feasible without incurring into breaking trade-offs;
while it may be regarded as a tiny step towards greener computing, it, along
with other initiatives, attempts to bring the end result closer to realization.

\section{Limitations And Future Work}
\label{sec:conclusions_limitations_and_future_works}

This section depicts some of the difficulties encountered, existing criticalities
and limitations, and relevant future ideas that can enhance and extend the
current architecture and implementation. Because one of the project's main points
is FOSS (Free and Open Source Software) compliance, anyone from anywhere can
contribute to it. The latter not only can increase the overall number of features
but also improve overall stability thanks to intrinsic testing in various use-case
scenarios and heterogeneous environments, which allows for the discovery and
correction of previously unknown bugs and/or errors. External contributions from
a wide variety of users are so valuable that they are frequently undervalued and
underestimated in current contexts. \\ %

\section{Internal Energy Source}
\label{subsec:conclusions_limitations_and_future_works_internal_energy_source}

The combination of the installation script and the external device for measuring
power consumption only conducts simple and basic readings on the instantaneously
drawn energy. This is ideal for desktop computers and servers that require an
external energy source, but it is incompatible with devices that use an internal
energy source (batteries), such as laptop computers. \\ %
If the existing implementation of the installation script is executed on a
laptop computer, the resulting power consumption readings are inaccurate since it
does not take into consideration the current absorbed energy from both the external
(plug) and internal (battery) sources, but only the first. As a result, the
obtained final power consumption data are always lower than what are in reality,
leading to general inefficiency and erroneous selection by the algorithm in upscaling
and downscaling operations, because the Server recognizes the node as a much
more power-efficient system than it is. \\ %
The script should determine whether the current node has internal energy sources
and, when conducting power consumption readings and calculations, adjust with a
more sophisticated but efficient and compatible algorithm that takes into
account for both external and internal absorbed energy. Moreover, further
testing on devices with an internal energy source is required to better
understand how the system is powered, as it can solely use the external source,
only use the internal source, or even employ both.

\section{Cluster Autoscaler Cloud Provider}
\label{subsec:conclusions_limitations_and_future_works_cluster_autoscaler_cloud_provider}

One of the project's most significant difficulties has been developing the
recluster Cloud Provider for the Cluster Autoscaler. There are three primary reasons
for this. \\ %
First off, the Cluster Autoscaler and the majority of the Kubernetes development
environment are built using the Go\footnote{\url{https://go.dev}} programming
language. I had never written a line of Go code before developing the \texttt{recluster}
Cloud Provider, so I had to learn the language from the beginning as well as the
package ecosystem, code style, and syntax. Although difficult at first, the effort
paid off in the end. \\ %
Second, neither official nor unofficial documentation exists that explains how
to create a Cloud Provider for the Cluster Autoscaler. My entire knowledge of developing
the \texttt{recluster} Cloud Provider is based on studying and evaluating the
code developed by the other Cloud Providers. Additionally, because each Cloud
Provider has a unique implementation, distinct Server logic, and a different API,
I had to determine the similarities and fundamental elements that almost all of them
share to transform it and make it compatible with reCluster. The lack of
documentation is understandable given that only a small number of entities
utilize and develop this technology, making it possible for them to access experts
and/or specialized consulting from the main CA developers themselves. Keep in
mind that the latter needs familiarity with the Go programming language, which was
covered in the preceding point. \\ %
Finally, it took longer than I anticipated to test the Cluster Autoscaler in
conjunction with the Server and the general architecture since almost every time
a new, minor error or bug that I hadn't considered appeared. However, the latter
was not only beneficial for reCluster; it additionally helped in identifying an
issue\footnote{\url{https://github.com/kubernetes/autoscaler/issues/5378}} in
the Cluster Autoscaler's main code, which was submitted to the maintainers and promptly
fixed within a week.

\section{Dashboard}
\label{subsec:conclusions_limitations_and_future_works_dashboard}

Currently, the only way for a user or administrator to interact with the cluster
is to use applications such as \texttt{curl}, \texttt{wget}, or similarities to directly
use the various set of GraphQL queries and mutation API, or use the GraphQL
Studio Explorer, which is not available in production and is only a nicer interface
for low-level interactions with the GraphQL API. As a result, there is a need to
create a UI dashboard, such as the one from Kubernetes\footnote{\url{https://kubernetes.io/docs/tasks/access-application-cluster/web-ui-dashboard}},
where all of the information provided by the Server is visually available without
involving the user in low-level GraphQL API interactions. The majority of
operations should be simple and intuitive, without needing the user to better understand
the requirements and outcomes. \\ %
Because the web application must consider the authentication mechanism, only
authorized and/or authenticated users may access specific information or do
certain operations. Furthermore, to provide a better overview of all the
information accessible to the user, the dashboard should be capable of generating
faster-to-understand graphs and/or charts. Consider a pie chart in which the
number of active/working nodes (green color) is visually compared against the
number of inactive nodes (gray color). The user can immediately understand the current
cluster capacity/demand, where if the chart is almost entirely green, there is an
overall high workload in the cluster, whereas if the chart is almost entirely
gray, there is an overall low workload in the cluster, and thus the number of
inactive nodes predominates. \\ %
With so many tools, technologies, and frameworks available for developing a web
application nowadays, the most difficult and critical phase is choosing which one
to use before even beginning development. \\ %

Also, there may be a desire to improve the simplicity with which each node is managed
without directly requiring \texttt{SSH} or similar technologies. \\ %
The latter can be accomplished by directly integrating a protected graphical
dashboard on each node that displays the various information of the associated
node as well as the ability to install, upgrade, or delete packages and perform
additional operations. \\ %
\texttt{YunoHost}\footnote{\url{https://yunohost.org}}, which leverages a simple
and easy-to-use, but incredibly powerful, web UI dashboard to completely
administer the node, is a starting point for what may be accomplished or developed.

\section{Mobile Devices}
\label{subsec:conclusions_limitations_and_future_works_mobile_devices}

Computers are the only systems that are currently used and supported in the
cluster. But, there are an enormous number of unused smartphone and tablet
devices that might be employed in the cluster. Most of these devices are ideal for
becoming worker nodes in a cluster because the previous owners replaced them due
to an upgrade to a newer model or because there is hardware damage, such as on
the screen or camera, but this has not affected important components such as the
main board. Furthermore, these devices have become so powerful in recent years that
they may be compared to the same performance as medium/low-end PCs while having a
significantly lower power consumption thanks to different internal architecture and
a prior design on energy consumption optimizations. \\ %
Developing for mobile devices might be challenging since there may be incompatibilities
between the software used for the cluster related to the internal Operating System
and missing libraries or Kernel requirements. A viable alternative is to
uninstall the existing operating system (\texttt{Android} or \texttt{iOS}) and
replace it with a pure GNU/Linux distribution designed exclusively for mobile
devices, such as \texttt{postmarketOS}\footnote{\url{https://postmarketos.org}}.
It should be noted that the overall compatibility of the Operating System with the
device varies from device to device, and the OS may need to be patched. Mobile development
should not be hard, however, this is the current status and/or only method for quickly
porting applications designed for GNU/Linux on PCs to mobile devices. \\ %
As previously stated, the power consumption of mobile devices is extremely small
compared to what they can provide. When compared to desktop computers and even laptop
computers, the performance-to-power-consumption ratio is enormous. There are no simple
techniques for remotely powering on a mobile device that does not involve custom
solutions. However, it should be mentioned that the sleep mode of mobile devices
is so efficient that a device in sleep mode consumes almost zero energy and may last
for more than a week on a single charge. As a result, if the remote wake-up solution
is too expensive or unfeasible, the concept of keeping a mobile device
constantly in sleep mode might be considered.

\section{Architectures}
\label{subsec:conclusions_limitations_and_future_works_architectures}

The cluster currently supports only \texttt{amd64} and \texttt{arm64} architectures.
Because the majority of the employed software and technologies support a wide
range of different architectures by default, such as \texttt{mips}\footnote{\url{https://wikipedia.org/wiki/MIPS_architecture}}
or \texttt{mipsel}, the remaining incompatible components, such as the Cluster Autoscaler,
can be ported and/or transformed to be compatible with additional and different
architectures. \\ %
Expanding the number of supported architectures allows more devices to be
incorporated into the cluster, giving them new life while also improving the cluster's
overall workload capacity.

\section{Node Selection Algorithm}
\label{subsec:conclusions_limitations_and_future_works_node_selection_algorithm}

The existing Node selection algorithm is effective, but it is a simplistic implementation
because it just depends on maximum and minimum power consumptions and, if both
are identical, falls back to performance data. \\ %
The existing algorithm works, despite its simplicity and limited accessible data
to perform all conceivable computations and evaluations. Nevertheless, there are
several aspects where it may be refined to accomplish a more suitable node's
final decision. As an example, while upscaling, the chosen node may have a lower
power consumption and possibly greater performance, whereas when downscaling, it
may have higher power consumption and potentially worse performance. Furthermore,
in future versions, the more efficient and sophisticated algorithm should take
into consideration other factors with a more advanced internal logic. There may be
nodes that are more efficient while under medium workload conditions, while
others are less efficient when under the same conditions but with lower minimum
and maximum power consumption coefficients. The existing approach is deterministic,
always preferring the second node type when upscaling and the first node type
when downscaling. The new algorithm should perform extra computations and evaluations
on which node type should bootstrap or terminate, and it's expected to follow
the current overall cluster state as well as its overall power consumption. \\ %
The algorithm may also be enhanced with a Kalman Filter\footnote{\url{https://wikipedia.org/wiki/Kalman_filter}}
for estimating the combination of power consumption and workload to
significantly reduce total (real) power consumption while also improving overall
Quality of Service (QoS). \\ %
The latter are merely possible ideas for what the new enhanced algorithm should
include and what possible advancements, such as lowering power consumption while
increasing overall performance, must be addressed.

\section{Performance Management}
\label{subsec:conclusions_limitations_and_future_works_performance_management}

\section{Scenarios Evaluation}
\label{subsec:conclusions_limitations_and_future_works_scenarios_evaluation}

% Scenario evaluation school, redundancy, disaster

\section{Server}
\label{subsec:conclusions_limitations_and_future_works_server}

% Server Simplification

% Better errors and node management

\begin{wrapfigure}
  {r}{.25\textwidth} %
  \centering
  \def\stackalignment{r}\stackunder{ \includegraphics[width=\linewidth]{images/logos/rust.png} } %
  {\scriptsize \parbox[t]{\linewidth}{ Source: \url{https://foundation.rust-lang.org/policies/logo-policy-and-media-guide}} }
  \caption{Rust logo}
\end{wrapfigure}

\section{Closing Remarks}
\label{sec:conclusions_closing_remarks}

Adapting existing tools that are specifically developed for big and energy-hungry
environments to a more sustainable and resource-aware design while maintaining
nearly the same capabilities for a more green future for the computing industry
is not only feasible, but also necessary.